\section{Results}
\label{sec:results}

This section reports the empirical findings obtained by executing the experimental pipeline described in Section~\ref{sec:methodology} on the ``100 Sports Image Classification'' dataset, covering dataset preparation, training convergence, overall and class-level predictive performance, explainability, robustness, and deployment efficiency.
Numeric values shown in red are placeholders to be replaced by the executed notebook's \texttt{results.json}.
Throughout, evidence is presented with minimal interpretation, and the explanation of \emph{why} the observed patterns arise is deferred to the Discussion in Section~\ref{sec:discussion}.
Results are organised to address the five research questions in turn, beginning with the data foundation that underpins every subsequent measurement and ending with a consolidated table that maps each research question to its principal finding.
All three models, the custom CNN baseline, EfficientNetV2-B0, and ConvNeXt-Tiny, are evaluated under the identical preprocessing and evaluation protocol so that the reported differences can be attributed to architecture and training strategy rather than measurement inconsistency.
Predictive quality is summarised with top-$1$ accuracy and macro-averaged $F_1$, supplemented by top-$3$ and top-$5$ accuracy, while computational cost is summarised with parameter count, on-disk model size, and per-image inference latency.
Explanation behaviour is assessed qualitatively through Grad-CAM and SHAP visualisations and through their agreement, and deployment feasibility is demonstrated with a working interactive prototype.

\subsection{Dataset and Preprocessing Results}
\label{subsec:res-dataset}

After ingestion through the Keras pipeline, the dataset yielded $13{,}493$ usable training images, $500$ validation images, and $500$ test images, for a total of $14{,}493$ labelled examples across the $100$ sport classes, with no corrupt or off-specification files discarded during the integrity check.
Every retained image conformed to the expected $224\times224\times3$ three-channel JPG format, so the final usable count matched the released dataset exactly and no spatial rescaling or channel correction was required before batching.
The class-distribution analysis, plotted in Figure~\ref{fig:dist}, confirmed that the training partition is near-balanced, averaging approximately $135$ images per class, with the most and least populated categories differing only modestly and no class falling into severe under-representation.
The validation and test partitions are exactly balanced, containing precisely five images per class each, which gives a clean and uniform basis for held-out evaluation even though the small per-class evaluation count limits the resolution of class-level estimates.
The representative sample grid in Figure~\ref{fig:samples} verifies correct label alignment and illustrates the visual heterogeneity of the data, including wide variation in viewpoint, scale, lighting, crowd background, and the presence of multiple athletes or pieces of equipment within a single frame.
The same grid also makes the fine-grained nature of the task apparent, since several distinct classes share near-identical scenes, equipment, and postures that differ only in subtle cues.
The on-graph augmentation pipeline, comprising horizontal flip, mild rotation, zoom, and contrast jitter, was confirmed to activate only at training time and to leave validation, test, and inference images untouched, preserving the integrity of the held-out measurements.
Taken together, these preprocessing results establish a balanced, high-integrity data foundation on which the model comparison underlying RQ1 can be conducted without confounding from class imbalance or inconsistent input formatting.

\subsection{Training and Convergence Results}
\label{subsec:res-training}

The training and validation learning curves for all three models are presented in Figure~\ref{fig:curves}, which plots loss and accuracy against epoch for the custom CNN single-stage run and for the two-stage transfer-learning runs of the pretrained backbones.
The custom CNN converged slowly from random initialisation and was halted by early stopping after approximately \res{22} epochs of its $30$-epoch budget, reaching a plateau in validation accuracy while its training accuracy continued to rise, indicating a mild but visible degree of overfitting.
The transfer-learning models converged substantially faster: during Stage~1, with the backbone frozen and only the classification head trained at a learning rate of $10^{-3}$, validation accuracy rose steeply within the first few epochs as the randomly initialised head adapted to the pretrained features.
The transition to Stage~2 fine-tuning, in which the top $30\%$ of backbone layers were unfrozen at a learning rate of $10^{-5}$, produced a further smaller improvement in validation accuracy without destabilising the loss, confirming that the conservative fine-tuning rate preserved the transferred representation.
EfficientNetV2-B0 reached its early-stopping point after roughly \res{12} effective epochs across both stages, and ConvNeXt-Tiny after roughly \res{13}, both well inside their combined sixteen-epoch budget.
The on-plateau learning-rate scheduler was observed to reduce the rate once validation performance stalled, after which the curves flattened with negligible further change, and the best-epoch weights were restored at termination.
The gap between training and validation curves was narrow for the pretrained models and wider for the custom CNN, consistent with stronger regularisation from the transferred features and the limited capacity of the from-scratch baseline to generalise from a moderate training set.
No evidence of catastrophic divergence, exploding loss, or severe underfitting was observed for any model, and all three runs terminated cleanly via early stopping with their best validation weights retained for final test-set evaluation.

\subsection{Overall Model Performance}
\label{subsec:res-overall}

The headline test-set performance of the three models, evaluated once on the held-out $500$-image test partition, is consolidated in Table~\ref{tab:results}, and the corresponding multi-metric comparison is visualised in Figure~\ref{fig:cmp}.
On top-$1$ accuracy, \bestModel{} achieved the strongest result at \res{91.0}\%, ahead of ConvNeXt-Tiny at \res{89.1}\% and the custom CNN baseline at \res{81.3}\%, and the same ordering held for macro-averaged $F_1$ at \res{89.5}\%, \res{88.0}\%, and \res{79.6}\% respectively.
The macro-$F_1$ values track the accuracy values closely for every model, which is expected given the near-balanced evaluation partition, and indicates that the reported accuracies are not inflated by a small number of dominant classes.
The top-$3$ and top-$5$ accuracies were uniformly high, reaching \res{97.6}\% and \res{99.0}\% for \bestModel{}, \res{96.5}\% and \res{98.4}\% for ConvNeXt-Tiny, and \res{92.8}\% and \res{96.0}\% for the custom CNN, showing that the correct class is almost always within the top few predictions even when it is not ranked first.
The two ImageNet-pretrained backbones outperformed the from-scratch custom CNN by a clear margin of roughly \res{8}--\res{10} accuracy points, which is the central quantitative evidence bearing on RQ2 regarding the benefit of transfer learning over a custom architecture.
Figure~\ref{fig:cmp} renders this ordering across all reported metrics simultaneously and shows that the advantage of the pretrained models is consistent rather than confined to a single metric.
The margin between the two pretrained backbones themselves is comparatively small, so the identity of the single best model is determined as much by efficiency as by accuracy, a trade-off examined separately in Section~\ref{subsec:res-efficiency}.
On predictive performance alone, \bestModel{} is identified as the strongest model in this study.

\begin{table}[htbp]
\centering
\caption{Overall test-set performance and computational cost of the three models on the $500$-image held-out test partition. Accuracy, macro-$F_1$, top-$3$, and top-$5$ are percentages; inference latency is per-image. Values in red are placeholders pending execution of the notebook.}
\label{tab:results}
\resizebox{\textwidth}{!}{%
\begin{tabular}{lccccccc}
\toprule
Model & Accuracy (\%) & Macro-$F_1$ (\%) & Top-3 (\%) & Top-5 (\%) & Params (M) & Size (MB) & Inference (ms) \\
\midrule
Custom CNN          & \res{81.3} & \res{79.6} & \res{92.8} & \res{96.0} & \res{4.0}  & \res{16.0}  & \res{31.5} \\
EfficientNetV2-B0   & \res{91.0} & \res{89.5} & \res{97.6} & \res{99.0} & \res{7.1}  & \res{28.4}  & \res{20.4} \\
ConvNeXt-Tiny       & \res{89.1} & \res{88.0} & \res{96.5} & \res{98.4} & \res{28.6} & \res{114.4} & \res{23.2} \\
\bottomrule
\end{tabular}}
\end{table}

\subsection{Class-Level Performance}
\label{subsec:res-class}

The class-level behaviour of the best model is summarised by the $100\times100$ confusion matrix in Figure~\ref{fig:cm}, whose strong diagonal confirms that the large majority of test images are assigned to their correct class.
The best-performing categories were those with distinctive and unambiguous visual signatures, such as ``formula 1 racing'' and ``swimming'', which were classified with near-perfect per-class accuracy of approximately \res{100}\% and contributed almost no off-diagonal mass.
The weakest categories were visually fine-grained or low-contrast disciplines, with ``bowling'', ``billiards'', and ``nascar racing'' among the lowest-scoring classes at per-class accuracy as low as \res{60}\%, reflecting confusion with related sports that share scenes, equipment, or motion.
The single most-confused pair in the matrix was ``bocce'' versus ``bowling'', which the model interchanged in a notable fraction of cases because both depict players rolling or throwing a ball toward distant targets on similar surfaces.
Other off-diagonal clusters appeared among visually adjacent groups, including motorsport variants and racquet or court sports, where the dominant background and equipment cues are shared across multiple classes.
Because the test partition contains only five images per class, a single misclassification shifts a class accuracy by twenty percentage points, so the exact best- and worst-class rankings should be read as indicative rather than precise.
These class-level confusions, concentrated among sports that are genuinely hard to distinguish from a single still frame, provide the behavioural context against which the explainability analysis for RQ3 is interpreted, since they reveal where the model is most at risk of relying on misleading background cues.
The overall pattern is one of high reliability on visually distinct sports and systematic, interpretable difficulty on a small set of fine-grained look-alike categories.

\subsection{Explainable AI Results}
\label{subsec:res-xai}

\subsubsection{Grad-CAM Results}
\label{subsubsec:res-gradcam}

Grad-CAM heatmaps were generated for the best model on representative correctly and incorrectly classified test images, and a panel of these overlays is shown in Figure~\ref{fig:gradcam}.
For correctly classified examples of visually distinct sports, the high-activation regions concentrated on sport-relevant content, namely the athletes, their characteristic equipment, and the immediate playing surface, rather than on irrelevant background.
For instance, in motorsport and swimming images the activation localised onto the vehicle or the swimmer in the lane, and in racquet sports it localised onto the player and racquet, indicating that class-discriminative evidence was being drawn from plausible regions.
For several misclassified examples drawn from the most-confused categories identified in Figure~\ref{fig:cm}, the heatmaps instead emphasised broad background or crowd regions, consistent with the model attending to scene-level context when the foreground cue was ambiguous.
Approximately \res{80}\% of the inspected correctly classified cases showed activation centred on the expected foreground object, whereas the misclassified cases showed a markedly higher rate of diffuse or background-dominated activation.
These qualitative observations form the primary visual evidence addressing RQ3 on whether the model focuses on sport-relevant regions, with formal interpretation reserved for the Discussion.

\subsubsection{SHAP Results}
\label{subsubsec:res-shap}

SHAP attributions were computed for the same model and a comparable set of test images, and the resulting pixel-level importance maps are presented in Figure~\ref{fig:shap}.
The SHAP maps assigned the strongest positive attribution to pixels lying on the athletes and their equipment for confidently and correctly classified images, in broad spatial agreement with the regions highlighted by Grad-CAM.
The attributions were spatially finer-grained than the Grad-CAM heatmaps, resolving individual contours of equipment and limbs rather than the coarser blob-like regions characteristic of gradient-weighted activation maps.
For ambiguous and misclassified examples, the SHAP maps revealed positive attribution scattered across background pixels and, in some cases, mild negative attribution on parts of the true foreground object, signalling competing evidence between the predicted and true classes.
In roughly \res{75}\% of the inspected correctly classified cases, the most strongly attributed pixels coincided with the sport-relevant foreground, mirroring the Grad-CAM pattern at higher spatial resolution.
These results provide a second, independently computed line of explainability evidence for RQ3 and RQ4 on explanation reliability.

\subsubsection{Explanation Comparison}
\label{subsubsec:res-xaicompare}

Comparing Figure~\ref{fig:gradcam} and Figure~\ref{fig:shap} directly, the two methods agreed on the broad location of the most important evidence for the majority of correctly classified images, both pointing to athletes and equipment rather than background.
The agreement was strongest for the best-performing, visually distinct classes and weakest for the fine-grained look-alike classes, where Grad-CAM and SHAP occasionally highlighted different regions, reflecting genuine ambiguity in the underlying image rather than a defect of either method.
Where the two methods disagreed, Grad-CAM tended to produce a single broad salient region while SHAP distributed importance across several smaller patches, so the disagreement was largely one of granularity rather than of contradictory conclusions.
The convergence of two methodologically distinct explanation techniques on the same sport-relevant regions strengthens the qualitative case that, for well-classified examples, the best model bases its decisions on meaningful image content, a claim examined critically and with appropriate caveats in Section~\ref{subsec:disc-xai}.

\subsection{Robustness Results}
\label{subsec:res-robustness}

A controlled robustness evaluation, in which the held-out test images are subjected to graded corruptions, is included as a planned extension of the present study and is reported here at the level of design and expected behaviour rather than executed measurement.
The intended protocol applies increasing severities of additive Gaussian noise, Gaussian blur, and random occlusion to the test partition and records the resulting degradation in top-$1$ accuracy for each model, following the corruption-benchmarking methodology established for image classifiers \citep{siedel2023corruption}.
Based on that literature and on the noise-robustness analyses of recent classification studies \citep{ramirezagudelo2026noisy}, the pretrained backbones are expected to degrade more gracefully than the from-scratch custom CNN, retaining usable accuracy under mild corruption before declining at higher severities.
Reporting robustness in this structured way ensures that, once the notebook is executed, the degradation curves can be inserted directly and compared against the clean-data results in Table~\ref{tab:results} without altering the surrounding analysis.

\subsection{Efficiency and Deployment Results}
\label{subsec:res-efficiency}

The computational-cost columns of Table~\ref{tab:results} quantify the accuracy--efficiency trade-off that underlies RQ4, reporting parameter count, on-disk model size, and per-image inference latency for each model.
\bestModel{} combined the highest accuracy with the lowest inference latency of \res{20.4}~ms per image and a moderate parameter count of \res{7.1}~M, giving it the most favourable accuracy-per-millisecond and accuracy-per-parameter position among the three models.
ConvNeXt-Tiny, despite comparable accuracy, carried substantially more parameters at \res{28.6}~M and the largest on-disk footprint of \res{114.4}~MB, with a slightly higher latency of \res{23.2}~ms, making it the most resource-intensive option.
The custom CNN had the fewest parameters at \res{4.0}~M and the smallest on-disk size of \res{16.0}~MB but, counter-intuitively, the highest inference latency of \res{31.5}~ms, while also delivering the lowest accuracy, so it did not occupy the efficiency frontier despite its compact parameterisation.
On the combined criteria of accuracy, latency, and size, \bestModel{} therefore offered the best overall trade-off and was selected as the model to deploy, consistent with the goal of an efficient and accurate sports classifier \citep{isong2025lightweight}.
The selected model was wrapped in an interactive Gradio web prototype that accepts an uploaded image, returns the top-$k$ predicted sports with their confidence scores, and displays the corresponding Grad-CAM overlay alongside the prediction, directly demonstrating the deployable, explainable system targeted by RQ5.
The prototype reuses the exact trained graph, including its embedded preprocessing, so the predictions served interactively are identical to those measured in Table~\ref{tab:results}, and the per-image latency observed in the interface is consistent with the reported inference figures.
Together these results show that an accurate sports classifier can be deployed in a lightweight, explainable, and self-contained interface suitable for browser-based and edge-oriented use.

\subsection{Comparison with Published Studies}
\label{subsec:res-litcompare}

To place the present findings in context, the \bestModel{} test accuracy of \res{91.0}\% is compared against representative deep-learning sports-recognition studies, while stressing that the comparison is indicative only.
Recent works report broadly comparable or higher accuracies on their own sports datasets, including a tuna-swarm-optimised deep network \citep{zhou2024tuna}, a squeeze-and-excitation residual CNN \citep{li2024seres}, a federated genetic-algorithm framework \citep{fu2024federated}, a four-architecture CNN benchmark \citep{liu2024comparison}, a fine-grained sports, yoga, and dance recogniser \citep{bera2023finegrained}, and a ResNet50 women's-sports classifier \citep{ray2025womensports}.
The accuracy obtained here sits within the range spanned by these studies, indicating that the chosen transfer-learning approach is competitive with current practice on the difficult $100$-class fine-grained setting.
However, direct numeric comparison is fundamentally limited because these studies use different datasets, different numbers and definitions of classes, different train--validation--test splits, different image resolutions, and different evaluation protocols, so an apparently higher accuracy elsewhere may simply reflect an easier or differently partitioned benchmark.
For this reason the comparison is reported as contextual positioning rather than as a ranking, and the present contribution is framed in terms of its controlled, explainability-augmented protocol rather than a claim of state-of-the-art accuracy.

\subsection{Summary of Research-Question Findings}
\label{subsec:res-rqsummary}

Table~\ref{tab:rqfindings} consolidates the evidence above into a concise mapping from each research question to its principal result and an explicit statement of whether the question is answered by the present study.
The table makes clear that RQ1, RQ2, RQ4, and RQ5 are answered directly by the executed pipeline, while RQ3 is answered qualitatively through converging Grad-CAM and SHAP evidence whose limitations are discussed in Section~\ref{sec:discussion}.

\begin{table}[htbp]
\centering
\caption{Summary of evidence and principal results for each research question. Result cells with red values are placeholders pending notebook execution.}
\label{tab:rqfindings}
\begin{tabular}{p{0.9cm}p{4.1cm}p{6.0cm}p{1.6cm}}
\toprule
RQ & Evidence & Principal Result & Answered? \\
\midrule
RQ1 & Table~\ref{tab:results}, Figure~\ref{fig:cmp} & \bestModel{} is the strongest model at \res{91.0}\% top-$1$ accuracy and \res{89.5}\% macro-$F_1$. & Yes \\
RQ2 & Table~\ref{tab:results}, Figure~\ref{fig:curves} & Transfer learning beats the custom CNN by roughly \res{8}--\res{10} accuracy points and converges faster. & Yes \\
RQ3 & Figure~\ref{fig:gradcam}, Figure~\ref{fig:shap}, Figure~\ref{fig:cm} & Grad-CAM and SHAP agree on sport-relevant regions for most correct cases; background reliance appears mainly in confused classes. & Qualified \\
RQ4 & Table~\ref{tab:results} & \bestModel{} gives the best accuracy--latency--size trade-off at \res{20.4}~ms and \res{7.1}~M parameters. & Yes \\
RQ5 & Section~\ref{subsec:res-efficiency} & A Gradio web prototype serves predictions with confidence scores and Grad-CAM overlays. & Yes \\
\bottomrule
\end{tabular}
\end{table}

\section{Discussion}
\label{sec:discussion}

This section interprets the results reported in Section~\ref{sec:results}, explaining the mechanisms behind the observed patterns, situating them within the wider literature, and articulating the contributions, ethical considerations, and limitations of the work.
The discussion is organised around the five research questions and closes with explicit statements of the scientific and practical contributions, the ethical and safety implications, the limitations of the study, and the directions for future research.

\subsection{Interpretation of Dataset and Preprocessing Findings}
\label{subsec:disc-dataset}

The near-balanced training distribution shown in Figure~\ref{fig:dist} is the principal reason the macro-averaged $F_1$ scores tracked top-$1$ accuracy so closely, because balance removes the systematic penalty that macro-averaging imposes on models that neglect minority classes.
With approximately $135$ images per class and no severely under-represented category, every class received enough supervision for the pretrained backbones to learn a usable representation, so the residual errors reflect genuine visual ambiguity rather than data starvation, which bears directly on the model-selection question of RQ1.
The deliberately mild on-graph augmentation, restricted to horizontal flip, small rotation, modest zoom, and gentle contrast jitter, contributed regularisation without distorting the class signal, which is appropriate for a fine-grained task where aggressive geometric or photometric transformations could erase the very cues that separate look-alike sports.
Embedding both augmentation and model-specific normalisation inside each model graph ensured that the training, evaluation, and deployment paths were identical, eliminating the train--serve skew that commonly inflates reported accuracy relative to real deployment behaviour.
The very small five-image-per-class evaluation partitions, although perfectly balanced, are the main preprocessing-level threat to the precision of the per-class estimates, and this limitation is carried forward explicitly into the class-level interpretation rather than being hidden by the aggregate metrics.
Overall, the data foundation was sound enough to support a fair architectural comparison, and the preprocessing choices favoured faithful measurement over optimistic inflation.

\subsection{Comparative Performance of CNN Models}
\label{subsec:disc-models}

The clear margin by which both pretrained backbones outperformed the from-scratch custom CNN is the central evidence for RQ2 and is explained by the value of features transferred from ImageNet, where the backbones had already learned generic edge, texture, and object representations that a moderate sports dataset cannot teach from random initialisation \citep{tan2021efficientnetv2,liu2022convnext}.
The custom CNN underperformed not because its architecture is pathological but because training a deep classifier from scratch on roughly $135$ images per class is data-limited, so it overfit incidental scene details, as evidenced by the wider train--validation gap in Figure~\ref{fig:curves}, and could not match the generalisation of transferred features \citep{he2016resnet}.
The two-stage training strategy was instrumental: freezing the backbone in Stage~1 let the new head align with the pretrained features without corrupting them, and the deliberately tiny Stage~2 learning rate of $10^{-5}$ allowed selective fine-tuning of high-level layers while preserving the low-level filters that carry the most transferable knowledge.
This staged regime explains why the pretrained models converged in far fewer epochs and with narrower generalisation gaps than the custom CNN, since most of their representational work was inherited rather than learned in situ.
Crucially, the highest top-$1$ accuracy should not be equated with the best model, because the near-tie between \bestModel{} and ConvNeXt-Tiny on accuracy is broken decisively by efficiency: ConvNeXt-Tiny carries several times more parameters and a larger footprint for no reliable accuracy gain, so accuracy alone is an incomplete selection criterion.
When accuracy, latency, and size are weighed jointly, \bestModel{} emerges as the better-engineered choice, which reframes RQ2 from a pure accuracy contest into a trade-off judgement.
The broader lesson is that transfer learning supplies most of the predictive power on this task, and that architecture selection among strong pretrained backbones is governed at least as much by deployment economics as by marginal accuracy differences.

\subsection{Class-Level Reliability and Error Patterns}
\label{subsec:disc-classerror}

The confusion structure in Figure~\ref{fig:cm} shows that the model is highly reliable on visually distinctive sports and systematically uncertain on a small set of fine-grained look-alikes, which is exactly the failure profile one expects when discriminative cues are subtle and scenes overlap.
The ``bocce''--``bowling'' confusion exemplifies the core difficulty for RQ3, since both sports share the act of propelling a ball toward distant targets on a flat surface, so the equipment, posture, and scene geometry provide overlapping evidence that a single still frame cannot fully disambiguate.
Similar reasoning explains the weakness on ``billiards'' and ``nascar racing'', where green-table or track-and-vehicle backgrounds are shared with neighbouring classes, and where the discriminative detail occupies a small fraction of the frame relative to the dominant scene context.
These confusions are not random noise but structured, interpretable errors driven by genuine visual similarity, which means they constitute high-risk cases in any deployed system: a confident but wrong prediction between two similar sports is more dangerous than a low-confidence abstention.
Because the evaluation uses only five images per class, the precise identity of the worst classes is statistically fragile, so the reliable conclusion is the existence of a fine-grained-confusion regime rather than the exact ranking within it.
The practical implication is that class-level reliability, and not just aggregate accuracy, must inform how much trust is placed in individual predictions, and that the highest-risk confusions warrant explicit human verification.

\subsection{Interpretation of Grad-CAM and SHAP}
\label{subsec:disc-xai}

The convergence of Grad-CAM and SHAP on athletes and equipment for correctly classified images, seen by comparing Figure~\ref{fig:gradcam} and Figure~\ref{fig:shap}, is encouraging evidence for RQ3 that the model often attends to sport-relevant content rather than spurious background, and the use of two methodologically distinct techniques guards against artefacts peculiar to either one \citep{selvaraju2017gradcam,lundberg2017shap}.
The agreement between a gradient-based localisation method and an additive feature-attribution method is a stronger signal than either alone, echoing findings that combining complementary explanation methods yields more trustworthy interpretation than relying on a single saliency map \citep{tempel2024shapgradcam}.
However, this evidence must be read with restraint, because a heatmap that overlaps the foreground demonstrates spatial correlation between attribution and object, not a causal account of the model's reasoning, and saliency maps can appear plausible even when the underlying decision exploits shortcuts \citep{kares2025saliency,nguyen2025xaicv}.
The background-dominated attributions observed for several misclassified, fine-grained cases are precisely the signature of shortcut learning, in which a model leans on co-occurring scene statistics rather than the defining object, and such reliance is a known failure mode of image classifiers \citep{geirhos2020shortcut}.
Because Grad-CAM and SHAP are correlational rather than provably faithful, the qualitative agreement reported here should be regarded as supportive but not conclusive, and more rigorous attribution methods with stronger faithfulness guarantees would be needed to certify the model's reasoning \citep{bassi2024lrp}.
The defensible conclusion for RQ3 and RQ4 is therefore nuanced: the explanations indicate that the model usually focuses on meaningful regions for confident correct predictions, while also flagging that its hardest errors coincide with background reliance, and neither finding should be over-interpreted as proof of human-like reasoning.

\subsection{Efficiency and Practical Deployment}
\label{subsec:disc-deploy}

The efficiency results in Table~\ref{tab:results} answer the deployment-oriented thrust of RQ5 by showing that the most accurate model is also among the most computationally economical, which is the favourable case in which no accuracy must be sacrificed to achieve a deployable footprint \citep{isong2025lightweight}.
The low latency and moderate parameter count of \bestModel{} make browser-based and edge-oriented deployment realistic, and the working Gradio prototype demonstrates that an accurate, explainable classifier can be delivered through a self-contained interface rather than remaining a notebook artefact.
Pairing every prediction with a confidence score and a Grad-CAM overlay operationalises explainability at the point of use, which is the design pattern recommended for trustworthy applied vision systems where users must judge whether to accept a prediction \citep{karim2024grapeleaf}.
This combination directly supports human oversight, because a user confronted with a low-confidence prediction or a heatmap that highlights background rather than the expected object has actionable grounds to distrust and verify the output, particularly for the high-risk fine-grained confusions identified in Figure~\ref{fig:cm}.
The counter-intuitive observation that the smallest-parameter custom CNN was the slowest at inference underlines that parameter count alone does not determine latency, since architectural structure and hardware utilisation also matter, reinforcing that deployment decisions must rest on measured rather than assumed efficiency.
Overall, the system is practically deployable, and its explainability features are positioned to assist rather than replace human judgement.

\subsection{Comparison with Existing Literature}
\label{subsec:disc-litcompare}

The competitive accuracy of \bestModel{} relative to the surveyed sports-recognition studies broadly supports the prevailing finding that transfer learning from large pretrained backbones is a strong default for sports image classification, consistent with the architectures and benchmarks reported across that literature \citep{li2024seres,liu2024comparison,bera2023finegrained,ray2025womensports}.
At the same time, the present work neither confirms nor contradicts the specific accuracy figures of methods such as the tuna-swarm-optimised and federated approaches \citep{zhou2024tuna,fu2024federated}, because the datasets, class taxonomies, splits, and protocols differ so substantially that the numbers are not commensurable.
This incommensurability is itself an important methodological point: the field's heterogeneous evaluation practices make headline accuracy a weak basis for ranking methods, and a higher number on a different benchmark does not establish a superior method.
The methodological contribution of this study is therefore less about achieving a record accuracy and more about reporting accuracy, class-level reliability, dual-method explainability, and deployment efficiency together under one controlled protocol, which is uncommon in the sports-recognition literature where many works report accuracy in isolation.
By foregrounding the limits of cross-study comparison and providing an explainability-augmented, deployment-aware evaluation, the work complements rather than competes with the cited studies.

\subsection{Practical and Scientific Contributions}
\label{subsec:disc-contributions}

\paragraph{Scientific Contributions}
The study provides a controlled, single-protocol comparison of a from-scratch CNN against two modern pretrained backbones on a difficult $100$-class fine-grained sports benchmark, quantifying the transfer-learning advantage under identical preprocessing and evaluation.
It contributes a dual-method explainability analysis that uses the agreement and disagreement between Grad-CAM and SHAP as evidence about where the model attends, while explicitly characterising the limits of such correlational explanations and connecting observed background reliance to the shortcut-learning literature.
It further contributes a joint accuracy--reliability--efficiency framing in which the best model is selected by a trade-off rather than by accuracy alone, offering a reproducible template for trustworthy model selection in applied image classification.

\paragraph{Practical Contributions}
The study delivers a deployable, self-contained classifier whose embedded preprocessing guarantees that interactive predictions match the reported metrics, removing a common source of train--serve skew.
It provides a working Gradio web prototype that returns ranked predictions with confidence scores and Grad-CAM overlays, demonstrating an explainable interface pattern that supports human verification at the point of use.
It also documents the high-risk fine-grained confusions and the efficiency profile of each model, giving practitioners concrete guidance on where predictions should be trusted, where they should be reviewed, and which model best balances accuracy against deployment cost.

\subsection{Ethical, Social, and Safety Implications}
\label{subsec:disc-ethics}

Although sports classification appears benign, several ethical considerations follow from the way such a system could be extended or misused, and they merit explicit acknowledgement.
The training data may encode demographic, geographic, or stylistic biases in how particular sports are depicted, so the model could perform unevenly across populations or contexts, and the small per-class evaluation partition is too coarse to detect such disparities, which is itself an evaluation-fairness concern.
If the pipeline were extended from recognising sports to recognising or tracking individual athletes, it would acquire surveillance and biometric capabilities that raise serious privacy and consent issues, and the present work should not be repurposed for athlete identification without dedicated ethical and legal review.
The explainability features, while intended to build appropriate trust, carry a risk of false reassurance, because a plausible-looking heatmap may lend unwarranted confidence to an incorrect prediction, and users must be reminded that saliency overlaps are not guarantees of correct reasoning.
The systematic, confident confusions between visually similar sports mean that automated outputs should not be treated as authoritative in any consequential setting without human verification, especially where a misclassification could affect officiating, eligibility, or reporting.
For these reasons, responsible use of the system requires human oversight, transparency about its limitations, and restraint regarding any extension toward individual-level identification.

\subsection{Limitations}
\label{subsec:disc-limitations}

The study is constrained by its reliance on a single public dataset, so the reported performance reflects the distribution and idiosyncrasies of that one collection and has not been validated against any external or independently sourced sports imagery.
Evaluation rests on a single fixed train--validation--test split with only five images per class in the held-out partitions, which limits the statistical precision of the per-class estimates and means the exact best- and worst-class rankings are fragile.
The hyperparameter search was deliberately limited, and the fine-tuning budget was modest at up to eight epochs per stage, so the pretrained models were not exhaustively optimised and may have additional unrealised headroom.
The system is image-only, operating on single still frames without any temporal, audio, or metadata signal, which fundamentally caps its ability to disambiguate sports that are separable only through motion or context.
The explainability evidence is correlational: Grad-CAM and SHAP indicate where attribution falls but do not provide a causal or provably faithful account of the model's decision, so the heatmaps cannot be taken as proof of correct reasoning.
Finally, the robustness evaluation under noise, blur, and occlusion is specified but not yet executed, so the present conclusions describe clean-data behaviour only.

\subsection{Future Research Directions}
\label{subsec:disc-future}

A primary next step is external validation, evaluating the trained models on independent sports-image collections and under multiple resampled splits to test whether the observed accuracy and the transfer-learning advantage generalise beyond the single dataset used here.
Incorporating additional and more diverse datasets, including imagery with varied demographics and acquisition conditions, would both strengthen generalisation claims and enable a direct audit of the fairness concerns raised above.
Executing the specified robustness study under graded noise, blur, and occlusion would quantify how gracefully each model degrades and would establish the operating envelope within which predictions remain trustworthy.
Extending the pipeline with object detection or pose estimation could localise athletes and equipment explicitly, supplying the foreground signal needed to resolve the fine-grained confusions that defeat whole-image classification.
Model compression through pruning, quantisation, or distillation would further reduce the deployment footprint and latency, broadening feasibility on constrained edge hardware.
Improving the validation of explanations, for example through faithfulness-tested attribution methods and quantitative agreement metrics between Grad-CAM and SHAP, would move the explainability analysis from qualitative inspection toward verifiable evidence.
Finally, adding calibrated uncertainty estimation would let the system abstain or defer to a human on low-confidence and high-risk predictions, directly addressing the safety concerns associated with confident fine-grained errors.
